\scanpage{8}
\uline{Cor.} Soit \deleted{\unsure{que}} $g\in G=O(E)$ transformant la calotte $(u,(x_i))$ en celle du lemme. Alors $g$ invarie l'icosaèdre $I=I_{uv}$

En effet, il est immédiat que l'icosaèdre associé à cette \inserted{deuxième} calotte est encore $I$.

Soit $\Gamma_I\inserted{{}=\Gamma_{u,v}}$ le sous-groupe de $G=O(E)$ qui laisse l'icosaèdre \inserted{$I=I_{u,v}$} invariant, $\Gamma'=\Gamma\cap G'$ le sous-groupe formé des rotations de $\Gamma$. Il est clair que $\Gamma'$ est d'indice 2 dans $\Gamma$ (N.B. $-1\deleted{\unsure{\notin}}\in\Gamma\cap\complement\Gamma'$) et que $\Gamma$ est d'où la produit
\[
\Gamma\simeq\mu_2\times\Gamma'
\]

\uline{Théorème} $\Gamma'$ opère de façon transitive sur l'ens. des sommets de l'icosaèdre, $I$, et de façon simplement transitive sur l'ens des couples $(x,\deleted{\unsure{y}}d)$ où $x$ est un sommet de l'icosaèdre, $I$ une arrête. $\Gamma$ opère de façon \inserted{simplement} transitive sur l'ens des triples $(x,y,z)$ de points de $I$ formant une face$^{*}$
% slanted
\margin{$*$ On peut aussi bien sur l'ens des triples d'arrêtes formant une face -- c'est évidemment équivalent}

On a oublié de préciser que par définition les \uline{arrêtes} de l'icosaèdre $I$ sont les \inserted{30} segments \deleted{de longueur $\ell$ qui joignent :} \inserted{suivants}
\[
\left\{\begin{aligned}
(u,x_i),&\qquad \deleted{\unsure{x_{i+1}}}(u',x'_i)\\
(x_i,x_{i+1}),&\qquad (x'_i,x'_{i+1})\\
(x_i,x'_{i-2})&\qquad (x'_i,x_{i-2})
\end{aligned}\right.\qquad i\in\mathbb Z/5\mathbb Z
\]
% slanted
\margin{Ce sont aussi les \unsure{ensembles} à deux éléments \inserted{distincts} $(x,y)$ de l'icosaèdre, de distance minimum (laquelle est $\ell$) -- sera plus facile à vérifier après transitivité de $\Gamma'$ sur l'ens des sommets}

De plus, on appelle \uline{faces} de l'icosaèdre une \deleted{partie} \inserted{des 20 \deleted{\unsure{ensembles}} parties} de $3$ \deleted{\unsure{éléments}} \inserted{trois} éléments \deleted{qui fait partie des \inserted{20} ensembles} suivants (ou les triangles qu'ils déterminent)
\[
\begin{aligned}
(u,x_i,x_{i+1})&\qquad (u',x'_i,x'_{i+1})\\
(x_i,x'_{i-2},\underset{\substack{\Vert\\x_{i+1}}}{x_{i\unsure{-}4}})&\qquad
(x'_i,x_{i-1},\underset{\substack{\Vert\\x'_{i+1}}}{x'_{i-4}})
\end{aligned}\qquad i\in\mathbb Z/5\mathbb Z
\]
% slanted
\margin{Ce sont aussi les triples \inserted{$(x,y,z)$} tels que $(x,y)$ $(y,z)$ et $(z,x)$ soient des arrêtes. Il y en a exactement \uline{cinq} issus de \deleted{sommet} \inserted{contenant} $u$ comme sommet. Ici aussi, ce sera immédiat à vérifier \uline{après} la transitivité.}

\uline{Démonstration} On sait déjà que les points $x_i$ \deleted{\unsure{sont situés}} sont congruents sous $\Gamma'$, et par symétrie les points $x'_i=-x_i$ aussi. De plus le lemme nous montre que $u$ est congruent à $x_0$, \deleted{$x_0$ est congru} (donc par symétrie $u'$ l'est à $x'_0$), et \deleted{par} $x_0$ \unsure{à} $x'_2$ -- donc tous les $x_i$ sont congruents. \deleted{Pour montrer} \inserted{Montrons} que $\Gamma'$ est \inserted{simplement} transitif sur l'ens des couples $(x,d)$ i.e des arrêtes \uline{orientées} $(x,y)$, \deleted{\unsure{il}} i.e que si $(x',y')$ en est une autre alors $\exists!\ g\in\Gamma'$, $gx'=x,gy'=y$. On peut supposer $x=u,y=x_0$, et par la transitivité déjà établie sur les sommets, $x'=u$. Donc \deleted{\unsure{en vertu comme}} \inserted{alors} $y'$ est un des $x_i$, \deleted{\unsure{et l'on sait}} \inserted{Il est clair} que le stabilisat\unsure{eur}
% archive sheet number: 7
